\ReportSection{5. Source Repository and Reproduction}
\ReportSubsection{5.1 Repository Structure}
Our source repository is
\url{https://github.com/james-helou/ECSE429-Project-B}.
We verified the executable code at commit \Path{e404151}. Later changes to
\Path{README.md} and \Path{BUG_REPORT.md} update the written explanations.
Section~8 gives the full commit ID and the execution records.

\begin{lstlisting}
ECSE429-Project-B/
+-- features/
|   +-- Story<N>*.feature         User stories and examples
+-- steps/
|   +-- Story<N>StepDefs.js       Steps for each story
|   +-- stepDefs.js              Shared steps and After hook
+-- TestUtil.js                  HTTP requests and record checks
+-- setup.js                     API startup, reset, and shutdown
+-- package.json                 Run commands and dependencies
+-- package-lock.json            Locked dependency versions
+-- runTodoManagerRestAPI-1.5.5.jar
+-- BUG_REPORT.md                Findings and reproduction steps
+-- README.md                    Setup and run instructions
\end{lstlisting}

\ReportSubsection{5.2 Environment and Run Instructions}
We ran the November 16 checks on Windows using Microsoft OpenJDK 17.0.12,
Node.js 22.13.1, npm 11.6.2, Cucumber.js 11.3.0, and Chai 4.5.0.
The suite requires Java and Node 18 or newer. From the repository directory,
install the dependencies and run the tests. On our Windows machine, we selected
the installed JDK with the following PowerShell commands:
\begin{lstlisting}
$env:JAVA_HOME = 'C:\Program Files\Microsoft\jdk-17.0.12.7-hotspot'
npm.cmd ci
npm.cmd run cucumber
\end{lstlisting}
The harness starts and stops the JAR automatically. Keep port 4567 free and run
one suite at a time. If an API instance is already answering on that port, the
harness shuts it down and its in-memory data is discarded.

\ReportSubsection{5.3 Changing the Order and Stopping the Service}
The default run uses random seed 765432. The commands below select a different
order and save its results, then run the service-stopped check:
\begin{lstlisting}
npx.cmd --no-install cucumber-js ./features -r ./steps `
  --order random:20261117 --format summary `
  --format json:results-seed-20261117.json
npm.cmd run cucumber:no-server
\end{lstlisting}
For the last command, leave the API stopped. Autostart is disabled, so every case
should fail at the readiness step. We expect exit code 1 because failures are
expected. To run only cases without the
\Path{@bug} tag, use \Path{npm run cucumber:no-bugs}.
